% GENERATED FILE. Edit canonical lessons, then run npm run generate:books.
% canonical-source-hash: fnv1a64:0934927e681d37b3
% canonical-lessons: ZH-C205-li, ZH-W205-li-apart, ZH-C205-gongli, ZH-W205-gong-public, ZH-C205-zuoyou, ZH-C205-chabuduo, ZH-C205-juli, ZH-R205-review-200-205

\chapter{How Far Is It?}
\label{ch:zh-how-far-is-it}
\hlchaptermodality{205}

\begin{chapteropening}
\textbf{By the end of this chapter:} \emph{I can say how far away a place is in kilometres, talk about distance and give a rough figure.}

Complete the last lesson of chapter 205: I can say how far away a place is in kilometres, talk about distance and give a rough figure.
\end{chapteropening}

\section[\texorpdfstring{\zh{离} (lí)}{离 (lí)}]{\zh{离} (lí) — from; away from (in distance)}

\label{lesson:ZH-C205-li}



\begin{quote}
Before the new one: say the Mandarin for video, then the Mandarin for to type on a keyboard.
\end{quote}

\subsection*{You'll want to know: \zh{离}}
\textbf{\zh{离}} — \emph{lí} — \textquotedblleft{}from; away from (in distance)\textquotedblright{}.

\begin{quote}
\emph{wǒ jiā lí xuéxiào hěn jìn} — My home is close to the school.
\end{quote}

The new character, in pieces: \zh{离} — ten strokes in three pieces: the dot and the bar that cap \zh{文}; two short strokes crossing inside a box open at the top; then a frame that hooks, holding a turning stroke and a dot. The next lesson puts it in your hand.

\begin{cousinweb}
\textbf{\zh{离}} (\emph{lí}) is inherited from Middle Chinese \emph{lje} (Guangyun, in Baxter's transcription), which says that being apart while near is \emph{lí}, and being apart while far is another word; in another tone it is glossed 'to leave'. Its use for distance from a place is not given there.
\end{cousinweb}

\begin{grammarlens}[title={What you've built}]
One more word for asking where things are.
\end{grammarlens}

\subsection*{Your turn}
\begin{itemize}
\raggedright
  \item \emph{Say it:} \emph{lí}
  \item \emph{Say it:} \emph{lí}, once more
  \item \emph{Say it:} \emph{wǒ jiā lí xuéxiào hěn jìn} — My home is close to the school.
  \item \emph{Recall:} say the Mandarin for video, then the Mandarin for to type on a keyboard, then say \emph{lí} again
\end{itemize}

\begin{tcolorbox}[breakable,colback=teal!4,colframe=teal!35!black,title={Before you move on}]
What does \zh{离} mean? (\textquotedblleft{}From; away from (in distance)\textquotedblright{}.) Say it once more.
\end{tcolorbox}
\section[\texorpdfstring{\zh{离} (lí)}{离 (lí)}]{\zh{离} — a cap, a box and a hooked frame}

\label{lesson:ZH-W205-li-apart}



\begin{quote}
The word first, then the character inside it.

\begin{itemize}
\raggedright
  \item \emph{Recall:} say \emph{lí} — \textbf{R1}, one lesson back, before it decays
\end{itemize}

\begin{quote}
\zh{文} \zh{开}
\end{quote}

\emph{Read it:} each one aloud, with nothing said first — \emph{wén}, \emph{kāi}
\end{quote}

\subsection*{Script — a cap, a box and a hooked frame}
\hlblockfigure{\detokenize{figures/ZH-W205-li-apart-filmstrip.pdf}}{How \zh{离} is written, stroke by stroke}

\begin{quote}
\zh{离}
\end{quote}

\textbf{Three pieces, top to bottom.}

\textbf{The top, two strokes:} the dot and the bar that cap \zh{文}.

\textbf{The middle, four strokes.} A short stroke falling to the left and a dot falling to the right, crossing like a small x. Then a box open at the top round them: a vertical down the left that turns along the foot, and a short vertical on the right.

\textbf{The bottom, four strokes.} A vertical down the left. A bar that turns down the right side and hooks to the left at its foot. Inside, a stroke that falls to the left and turns right, then a dot.

2 + 4 + 4 = ten.

\subsection*{Your turn}
\begin{itemize}
\raggedright
  \item \emph{Write it:} \zh{离}, all ten, with the model in front of you
  \item \emph{Cover:} write it again from memory, and count the pieces, not the strokes
\end{itemize}

\begin{tcolorbox}[breakable,colback=teal!4,colframe=teal!35!black,title={Before you move on}]
How many strokes in \zh{离}? (\textbf{Ten}.) What caps it? (\textbf{The dot and the bar of \zh{文}}.) Which stroke hooks? (\textbf{The bar that turns down the right side}.)

Source: \href{https://raw.githubusercontent.com/chanind/hanzi-writer-data/68d10a4b21150cae5e1ebbd223eed289cf32d90c/data/\%E7\%A6\%BB.json}{Hanzi Writer Data \zh{离}, pinned stroke paths and medians}.
\end{tcolorbox}
\section[\texorpdfstring{\zh{公里} (gōnglǐ)}{公里 (gōnglǐ)}]{\zh{公里} (gōnglǐ) — kilometre}

\label{lesson:ZH-C205-gongli}



\begin{quote}
Before the new one: say the Mandarin for to type on a keyboard, then the Mandarin for away from.
\end{quote}

\subsection*{You'll want to know: \zh{公里}}
\textbf{\zh{公里}} — \emph{gōnglǐ} — \textquotedblleft{}kilometre\textquotedblright{}.

\begin{quote}
\textbf{\zh{书店在东边离这儿两公里}} \emph{shūdiàn zài dōngbian lí zhèr liǎng gōnglǐ} — The bookshop is to the east, two kilometres from here.
\end{quote}

The new character, in pieces: \zh{公} — four strokes: \zh{八} on top, then a stroke that falls to the left and turns right, and a dot. The next lesson puts it in your hand.

\begin{cousinweb}
\textbf{\zh{公里}} (\emph{gōnglǐ}) is built from \zh{公} 'common, public' and \zh{里} \emph{lǐ}. The first is inherited from Middle Chinese \emph{kuwng} (Guangyun, in Baxter's transcription), glossed 'common' and 'without bias'; the dictionary glosses \zh{里} as a neighbourhood of twenty-five households. Its use as a measure of distance, and the pair's use for the kilometre, are not given there, so they are not settled.
\end{cousinweb}

\begin{grammarlens}[title={What you've built}]
One more word for asking where things are.
\end{grammarlens}

\subsection*{Your turn}
\begin{itemize}
\raggedright
  \item \emph{Say it:} \emph{gōnglǐ}
  \item \emph{Say it:} \emph{gōnglǐ}, once more
  \item \emph{Say it:} \emph{shūdiàn zài dōngbian lí zhèr liǎng gōnglǐ} — The bookshop is to the east, two kilometres from here.
  \item \emph{Recall:} say the Mandarin for to type on a keyboard, then the Mandarin for away from, then say \emph{gōnglǐ} again
\end{itemize}

\begin{tcolorbox}[breakable,colback=teal!4,colframe=teal!35!black,title={Before you move on}]
What does \zh{公里} mean? (\textquotedblleft{}Kilometre\textquotedblright{}.) Say it once more.
\end{tcolorbox}
\section[\texorpdfstring{\zh{公} (gōng)}{公 (gōng)}]{\zh{公} — \zh{八} over a turn and a dot}

\label{lesson:ZH-W205-gong-public}



\begin{quote}
The word first, then the character inside it.

\begin{itemize}
\raggedright
  \item \emph{Recall:} say \emph{gōnglǐ} — \textbf{R1}, one lesson back, before it decays
\end{itemize}

\begin{quote}
\zh{八} \zh{里}
\end{quote}

\emph{Read it:} each one aloud, with nothing said first — \emph{bā}, \emph{lǐ}
\end{quote}

\subsection*{Script — \zh{八} over a turn and a dot}
\hlblockfigure{\detokenize{figures/ZH-W205-gong-public-filmstrip.pdf}}{How \zh{公} is written, stroke by stroke}

\begin{quote}
\zh{公}
\end{quote}

\textbf{Two pieces, top before bottom.}

\textbf{The top piece, \zh{八}, two strokes:} a stroke falling to the left, then a longer one falling to the right, spread wide apart.

\textbf{The bottom piece, two strokes.} A stroke that falls to the left and turns to run right along the foot, then a dot at its right end. The same pair sits in the top left corner of \zh{能}.

2 + 2 = four.

\subsection*{Your turn}
\begin{itemize}
\raggedright
  \item \emph{Write it:} \zh{公}, all four, with the model in front of you
  \item \emph{Cover:} write it again from memory, and count the pieces, not the strokes
\end{itemize}

\begin{tcolorbox}[breakable,colback=teal!4,colframe=teal!35!black,title={Before you move on}]
How many strokes in \zh{公}? (\textbf{Four}.) What is the top piece? (\textbf{\zh{八}}.) Where has your hand made the bottom piece before? (\textbf{In the top left of \zh{能}}.)

Source: \href{https://raw.githubusercontent.com/chanind/hanzi-writer-data/68d10a4b21150cae5e1ebbd223eed289cf32d90c/data/\%E5\%85\%AC.json}{Hanzi Writer Data \zh{公}, pinned stroke paths and medians}.
\end{tcolorbox}
\section[zuǒyòu]{zuǒyòu — about; or so (after a number)}

\label{lesson:ZH-C205-zuoyou}



\begin{quote}
Before the new one: say the Mandarin for away from, then the Mandarin for kilometre.
\end{quote}

\subsection*{You'll want to know: zuǒyòu}
\emph{zuǒyòu} — \textquotedblleft{}about; or so (after a number)\textquotedblright{}.

\begin{quote}
\emph{zǒulù èrshí fēnzhōng zuǒyòu} — It's about twenty minutes on foot.
\end{quote}

Its characters come later; for now it is a word for the ear.

\begin{cousinweb}
\emph{zuǒyòu} is built from \emph{zuǒ} 'left' and \emph{yòu} 'right', both inherited from Middle Chinese, \emph{tsaX} and \emph{hjuwX} (Guangyun, in Baxter's transcription): a little to either side of the number.
\end{cousinweb}

\begin{grammarlens}[title={What you've built}]
One more word for asking where things are.
\end{grammarlens}

\subsection*{Your turn}
\begin{itemize}
\raggedright
  \item \emph{Say it:} \emph{zuǒyòu}
  \item \emph{Say it:} \emph{zuǒyòu}, once more
  \item \emph{Say it:} \emph{zǒulù èrshí fēnzhōng zuǒyòu} — It's about twenty minutes on foot.
  \item \emph{Recall:} say the Mandarin for away from, then the Mandarin for kilometre, then say \emph{zuǒyòu} again
\end{itemize}

\begin{tcolorbox}[breakable,colback=teal!4,colframe=teal!35!black,title={Before you move on}]
What does zuǒyòu mean? (\textquotedblleft{}About; or so (after a number)\textquotedblright{}.) Say it once more.
\end{tcolorbox}
\section[chàbuduō]{chàbuduō — about; almost}

\label{lesson:ZH-C205-chabuduo}



\begin{quote}
Before the new one: say the Mandarin for kilometre, then the Mandarin for or so.
\end{quote}

\subsection*{You'll want to know: chàbuduō}
\emph{chàbuduō} — \textquotedblleft{}about; almost\textquotedblright{}.

\begin{quote}
\emph{cóng zhèr dào chēzhàn chàbuduō liǎng gōnglǐ} — From here to the station is almost two kilometres.
\end{quote}

Its characters come later; for now it is a word for the ear.

\zh{不} loses its tone in the middle of this word: chàbuduō.

\begin{cousinweb}
\emph{chàbuduō} is built from \emph{chà} 'short of', \zh{不} 'not' and \emph{duō} 'much', all inherited from Middle Chinese, \emph{tsrhae}, \emph{pjut} and \emph{ta} (Guangyun, in Baxter's transcription): not much short.
\end{cousinweb}

\begin{grammarlens}[title={What you've built}]
One more word for asking where things are.
\end{grammarlens}

\subsection*{Your turn}
\begin{itemize}
\raggedright
  \item \emph{Say it:} \emph{chàbuduō}
  \item \emph{Say it:} \emph{chàbuduō}, once more
  \item \emph{Say it:} \emph{cóng zhèr dào chēzhàn chàbuduō liǎng gōnglǐ} — From here to the station is almost two kilometres.
  \item \emph{Recall:} say the Mandarin for kilometre, then the Mandarin for or so, then say \emph{chàbuduō} again
\end{itemize}

\begin{tcolorbox}[breakable,colback=teal!4,colframe=teal!35!black,title={Before you move on}]
What does chàbuduō mean? (\textquotedblleft{}About; almost\textquotedblright{}.) Say it once more.
\end{tcolorbox}
\section[jùlí]{jùlí — distance}

\label{lesson:ZH-C205-juli}



\begin{quote}
Before the new one: say the Mandarin for or so, then the Mandarin for almost.
\end{quote}

\subsection*{You'll want to know: jùlí}
\emph{jùlí} — \textquotedblleft{}distance\textquotedblright{}.

\begin{quote}
\emph{liǎng ge chēzhàn de jùlí shì yì gōnglǐ} — The two stations are a kilometre apart.
\end{quote}

Its characters come later; for now it is a word for the ear.

\begin{cousinweb}
\emph{jùlí} is built from \emph{jù} and \zh{离} 'away from', both inherited from Middle Chinese, \emph{gjoX} and \emph{lje} (Guangyun, in Baxter's transcription). The dictionary glosses the first as the spur on a cock's leg; how it came to mean being apart is not settled.
\end{cousinweb}

\begin{grammarlens}[title={What you've built}]
One more word for asking where things are.
\end{grammarlens}

\subsection*{Your turn}
\begin{itemize}
\raggedright
  \item \emph{Say it:} \emph{jùlí}
  \item \emph{Say it:} \emph{jùlí}, once more
  \item \emph{Say it:} \emph{liǎng ge chēzhàn de jùlí shì yì gōnglǐ} — The two stations are a kilometre apart.
  \item \emph{Recall:} say the Mandarin for or so, then the Mandarin for almost, then say \emph{jùlí} again
\end{itemize}

\begin{tcolorbox}[breakable,colback=teal!4,colframe=teal!35!black,title={Before you move on}]
What does jùlí mean? (\textquotedblleft{}Distance\textquotedblright{}.) Say it once more.
\end{tcolorbox}
\section[(dialogue)]{Looking back: from Booking a Room to How Far Is It?}

\label{lesson:ZH-R205-review-200-205}



\begin{quote}
Say the Mandarin for hotel and to book ahead.
\end{quote}

\subsection*{Your turn}
\begin{itemize}
\raggedright
  \item \emph{Say it:} hotel, to book ahead, single room, to check in, to check out
  \item \emph{Say it:} postage stamp, courier service, email, to send electronically, to get something sent
  \item \emph{Say it:} to scan a code, receipt, to take goods back, customer, checkout
  \item \emph{Say it:} to rent, landlord, rent for a home, a deposit paid in advance, contract
  \item \emph{Say it:} website, to log in, to click, video, to type on a keyboard
  \item \emph{Say it:} away from, kilometre, or so, almost, distance
  \item \emph{Read it:} \zh{东} \zh{离} \zh{公} — each one aloud, then a word that holds it
\end{itemize}

\begin{tcolorbox}[breakable,colback=teal!4,colframe=teal!35!black,title={Before you move on}]
Hotel? (\emph{bīnguǎn}.) To book ahead? (\emph{yùdìng}.) Single room? (\emph{dānrénjiān}.) To check in? (\emph{rùzhù}.) To check out? (\emph{tuìfáng}.) Postage stamp? (\emph{yóupiào}.) Courier service? (\emph{kuàidì}.) Email? (\emph{yóujiàn}.) To send electronically? (\emph{fāsòng}.) To get something sent? (\emph{shōudào}.) To scan a code? (\emph{sǎomǎ}.) Receipt? (\emph{shōujù}.) To take goods back? (\emph{tuìhuò}.) Customer? (\emph{gùkè}.) Checkout? (\emph{shōuyíntái}.) To rent? (\emph{zū}.) Landlord? (\textbf{\zh{房东}}, \emph{fángdōng}.) Rent for a home? (\emph{fángzū}.) A deposit paid in advance? (\emph{yājīn}.) Contract? (\emph{hétong}.) Website? (\textbf{\zh{网站}}, \emph{wǎngzhàn}.) To log in? (\emph{dēnglù}.) To click? (\emph{diǎnjī}.) Video? (\emph{shìpín}.) To type on a keyboard? (\textbf{\zh{打字}}, \emph{dǎzì}.) Away from? (\textbf{\zh{离}}, \emph{lí}.) Kilometre? (\textbf{\zh{公里}}, \emph{gōnglǐ}.) Or so? (\emph{zuǒyòu}.) Almost? (\emph{chàbuduō}.) Distance? (\emph{jùlí}.)
\end{tcolorbox}
